% DRAFT contribution to the Phase 2 related-work section (plan P1 / Rubric: "DL literature review, 4 DL papers").
% Written by Bhavya for Priyansh to merge or reword. Every factual claim below was checked against the
% paper's arXiv / venue listing; the sentences after "In our work" are our own framing, not the papers' claims.
% Keys: references_gnn.bib (kipf2017gcn, xu2019gin, hu2020pretrain, li2018oversmooth, oono2020expressive)
% and ../phase1/references.bib (gilmer2017mpnn, xiong2020attentivefp, yang2019analyzing).
\subsection{Graph neural networks for molecules}
\textbf{Graph convolutional networks.} Kipf and Welling~\cite{kipf2017gcn} derive a convolution layer for graph-structured data from a first-order approximation of spectral graph convolutions; its cost grows linearly with the number of edges, and its hidden layers encode both graph structure and node features. In our work it is the simplest message-passing baseline: it averages neighbours with a degree normalisation and has no input for bond types.

\textbf{How expressive are GNNs?} Xu et al.~\cite{xu2019gin} build a theoretical framework for the discriminative power of message-passing networks. They show that popular variants such as GCN and GraphSAGE fail to distinguish some simple graph structures, and propose the Graph Isomorphism Network, whose discriminative power matches the Weisfeiler--Lehman graph isomorphism test. In our work this motivates summing (not averaging) neighbour messages in the bond-aware model.

\textbf{Edge-aware layers and pre-training for molecules.} Hu et al.~\cite{hu2020pretrain} study pre-training GNNs at both node and whole-graph level and report gains of up to 9.4\% absolute ROC-AUC over non-pre-trained models on molecular property and protein-function prediction, while pre-training at only one level gave limited gains and sometimes negative transfer. The edge-feature variant of the GIN layer used in that paper is the layer implemented as \texttt{GINEConv} in PyTorch Geometric, which we use to feed bond features into the messages. We do not use pre-training.

\textbf{Depth and over-smoothing.} Li et al.~\cite{li2018oversmooth} show that the GCN convolution is a special form of Laplacian smoothing, which explains why it works but also why stacking many layers can make vertex features indistinguishable. Oono and Suzuki~\cite{oono2020expressive} analyse GCNs as dynamical systems and prove that, under conditions on the weights, the output of a deep network loses all information except the connected components and node degrees as depth grows. These are asymptotic results on citation-style benchmark graphs; whether depth hurts molecular graphs of median diameter nine bonds at 2--5 layers is tested empirically by our layer ablation.

\textbf{Position in the molecular literature.} Message passing was unified for chemistry by Gilmer et al.~\cite{gilmer2017mpnn}, and AttentiveFP~\cite{xiong2020attentivefp} adds attention at the atom and molecule level. Yang et al.~\cite{yang2019analyzing} caution that well-tuned fingerprint baselines are often competitive with learned representations, so we report GNNs against the classical models on identical splits instead of assuming they win.
